\documentclass[11pt]{article}

\usepackage[T1]{fontenc}
\usepackage[margin=1in]{geometry}
\usepackage{amsmath,amssymb,mathtools}

\newcommand{\Failed}{\mathsf{fail}}
\newcommand{\Value}{\mathsf{ok}}
\newcommand{\Sat}{\mathsf{accept}}
\newcommand{\Reject}{\mathsf{reject}}
\newcommand{\owner}{\mathsf{owner}}
\newcommand{\deps}{\mathsf{deps}}
\newcommand{\subject}{\mathsf{subject}}
\newcommand{\lean}[1]{\texttt{\detokenize{#1}}}

\title{Deferred Relational Guards with Null-Vacuous Failure}
\author{}
\date{}

\begin{document}
\maketitle

\paragraph{Scope and status.}
This note specifies deferred relational guards and their failure-aware source
execution, as formalized in \lean{Vegas.Source}. A guard is declared at the
commitment it constrains and checked later, when its inputs are public. The
deferral is syntactic: the program text fixes which reveal publishes each
private cell, hence which reveal completes each guard, so a guard is checked
exactly once, with a binary verdict, and no publication carries a runtime
status. The compiler consumes this \lean{SourceProgram} semantics and has
checked exact outcome and unilateral-deviation laws through the ordered graph
to its ideal pending-message runtime. This note isolates the source
publication mathematics; the separate compiler certificates carry the runtime
information and service obligations.

\section{Heterogeneous publications}
\label{sec:publications}

Let $I$ be a finite set of commitments. Persistent private inputs are outside
$I$: they require no publication and cannot occur in a guard's support.
Each commitment $i:I$ has an ordinary
value type $A_i$ and an owner $\owner(i)$.  Its binding is
\[
  B_i \in \mathsf{Result}(A_i)
      := \{\Failed\}\;\mathbin{+}\;\{\Value(a):a\in A_i\},
\]
where $\Value(a)$ is an ordinary candidate and $\Failed$ an unopenable one.
The binding is fixed at commitment and never changes.

A cell becomes public at its reveal, which creates a fresh public cell holding
a $\mathsf{Result}(A_i)$: either the ordinary value that was disclosed and
accepted, or $\Failed$, a public null outside $A_i$.  In particular, failure
does not fabricate an inhabitant of an ordinary domain, and it does not change
the raw binding.  This distinction is essential for programs which retain and
later use an original private value.

Public cells are immutable: a reveal writes one new cell and no transition
rewrites an earlier one.  Hence a later guard failure never retracts an earlier
public value.  Where a private cell's result is public is itself static: the
syntax determines, for each cell, the public cell its reveal creates.  We write
$P$ for the resulting public state and $P_i$ for the published result of cell
$i$, defined once $i$ has been revealed.

\section{Relational obligations}

A publication guard $g$ consists of
\begin{itemize}
\item a subject cell $s_g$;
\item a finite support $D_g\subseteq I$, with $s_g\in D_g$; and
\item a decidable relation
\[
  R_g : \prod_{i\in D_g} A_i \longrightarrow \mathsf{Bool}
\]
on ordinary typed values.
\end{itemize}
The subject records which guarded source choice introduced the obligation.
Requiring $s_g\in D_g$ prevents a guard from becoming false before its own
publication exists; it also handles unary and constant-false guards uniformly.

Because reveals are syntactic, the last member of $D_g$ to be revealed is fixed
by the program text.  Call that reveal the one which \emph{completes} $g$: the
reveal after which every member of $D_g$ is public and before which some member
was not.  Every declared guard has exactly one completing reveal, and its
verdict there is binary:
\[
 \mathsf{check}_g(P)=
 \begin{cases}
   \Sat, & \text{if some $i\in D_g$ has $P_i=\Failed$},\\
   \Sat, & \text{if all are ordinary values and $R_g$ is true},\\
   \Reject, & \text{if all are ordinary values and $R_g$ is false}.
 \end{cases}
\]
Thus failure is null-vacuous: a relation is evaluated only on an entirely
ordinary tuple, and any failed member of its static support discharges it.
No case reads an unopened private binding, because at the completing reveal
every member of $D_g$ is already public.  Static support is semantic, including
dependencies in dead branches, so ordinary Boolean equivalence alone does not
permit a support-changing rewrite: it can move the completing reveal.

A state is \emph{consistent} when every guard whose support is entirely public
accepts.  A guard whose support is not yet public is a real deferred
obligation, neither evidence that its relation is true nor a violation.

\section{Atomic resolution}

Suppose cell $i$ has binding $B_i$ and is revealed, with disclosure decision
$d\in\mathsf{Bool}$.  The proposal is
\[
 c =
 \begin{cases}
   B_i,&d=\mathsf{true},\\
   \Failed,&d=\mathsf{false},
 \end{cases}
\]
so withholding and an unopenable binding propose the same thing.  Let $G_i$ be
the guards this reveal completes, and $P^{c}$ the public state extended by
$P_i=c$.  If every $g\in G_i$ has $\mathsf{check}_g(P^c)=\Sat$, publish $c$;
otherwise publish $\Failed$.  No other cell changes.  In an implementation, the
tentative check and final write are one atomic application transition.

Candidate evaluation at this transition is public-only.  The relation receives
ordinary values only from public cells and from the current proposal; no
private binding is inspected.  Binding-to-opening authentication is a separate
check which supplies the candidate for the current cell.

An invalid opening and a missing opening have the same source result only when
the enclosing scheduler invokes this resolution at the same declared reveal or
finalization point.  Earlier malformed or mismatching target attempts must be
source stutters.  They may still be observable in a mempool or message runner;
erasing that extra transport information is a later strategic proof obligation,
not a property of this publication state.

In particular, an authentic opening whose ordinary value violates a guard may
still expose that value in its transaction payload. Recording $\Failed$ as the
source publication does not erase the observed payload or keep it secret.
Those observations must remain in the target game and its deviation argument.

\subsection{Constructive consistency invariant}
\label{sec:invariant}

The central lemma is local.

\paragraph{Preservation lemma.}
If $P$ is consistent and cell $i$ has not been revealed, then revealing $i$ by
the rule above produces a consistent state $P'$ and adds exactly one public
cell.

\paragraph{Proof sketch.}
Consider any guard $g$ whose support is public in $P'$.  If $g\notin G_i$, its
support was already public in $P$, its values are unchanged, and consistency of
$P$ gives $\mathsf{check}_g(P')=\Sat$.  If $g\in G_i$ and the proposal was
published, then $P'=P^c$ and $g$ accepted by the acceptance test.  If $g\in
G_i$ and the proposal was rejected, then $P'_i=\Failed$ with $i\in D_g$, so
null-vacuity makes $g$ accept.  This covers all guards.  Notice that the proof
never edits the subject of a guard retroactively: it fails the current
publication which exposed the conflict.

Induction over the declared source reveal order gives consistency of every
reachable state.  Since the syntax reveals every private cell, at return every
guard's support is public, so every guard accepts.

This theorem is not an arbitrary-order confluence result.  Reveal order can
change which cell completes a guard and therefore which cell fails.  Source
reveal order is part of the semantics and must be retained by a compiler.

\section{Ownership and failure attribution}

Consistency alone says nothing about whose deviation caused failure.  The
needed local ownership condition is
\[
  \text{$i$ not yet revealed}\ \wedge\ i\in D_g
  \quad\Longrightarrow\quad
  \owner(i)=\owner(s_g).
\]
Call this \emph{owned unrevealed support}.  A guard rejected at the reveal of
$i$ satisfies $i\in D_g$, since only the guards this reveal completes are
checked and every such guard reads $i$.  Owned unrevealed support therefore
shows that the failed current publisher and the subject author have the same
owner.  If several guards reject simultaneously, their subjects have that same
owner as well.

This is an equality of cell owners, not yet a proof of causal blame in a
message runtime. The source language has no caller-authentication or deadline
service interface. A host must enforce who may bind, open, or decline a site's
publication, and establish when failure can be attributed to that owner.
In particular, an outsider's invalid message must not immediately fail someone
else's publication. Neither a timeout without timely service nor a reverted
call establishes intentional withholding. Voluntary failure and unopenable
bindings need this host-level attribution too; they need not falsify any guard.

The condition matches the information discipline of the source language.  When
a player makes a guarded choice, public dependencies in that player's view have
already been published.  A dependency which remains sealed and is readable by
the chooser is the chooser's own initial or committed value.
This provenance should be proved from source visibility and operation order.
It must not be replaced by a blanket assertion that all arbitrary relational
guards have good attribution.

Without the condition, the semantics is operationally consistent but not a
unilateral quitting model.  For example, let player $p$ author a relation
$R(y,x)$ whose still-hidden $x$ is owned by player $q$.  Publish $y$ first.
Player $p$ can choose it so that $q$'s later publication of $x$ closes the
relation falsely, causing $q$'s publication to fail.  The failure cannot honestly be
charged to $p$'s source quit while the public state records $q$ as the last
publisher, nor to $q$ if $q$ never had the relation in its information set.

Ownership may optionally be made persistent: after one owned failure, later
publications controlled by that owner can be forced to fail.  Such a rule is
monotone and preserves the consistency proof, but it changes later source
choices and must be specified as an additional source transition.  It is not a
consequence of null-vacuity alone.

\section{Examples that distinguish the semantics}

\subsection{Reverse disclosure}

Let one player bind $x,y:\mathsf{Bool}$ with obligation $y=x$, then reveal $y$
before $x$.  The reveal of $y$ completes no guard, so publishing $y=b$ leaves
the obligation undecided.  The reveal of $x$ completes it, and there,
\begin{itemize}
\item $x=b$ publishes an ordinary value and satisfies the obligation;
\item $x\ne b$ makes the tentative tuple reject, so the final publication of
      $x$ is failed; and
\item no valid opening for $x$ produces the same failed publication at that
      reveal's finalization point.
\end{itemize}
The earlier public $y=b$ is never retracted.  The final lifted relation is true
vacuously in the failure branches.

Before the completing reveal, the public value has its ordinary payload type
but does not carry a certificate that $y=x$. A program may use the disclosed data,
but a payment justified by the relation must account for the outstanding
obligation. For example, settlement can wait for both publications and
penalize the responsible owner's failure. Irrevocably paying for an unchecked
claim without recoverable collateral would not be made safe by this semantics.

Reversing the reveal order reverses which mismatching cell fails.  Both results
are consistent.  They are different source outcomes, so a compiler may not
reorder these reveals under a confluence argument.

\subsection{Heterogeneous relation}

Let $x:\mathsf{Bool}$ and $n:\mathsf{Fin}(3)$ be owned by one player, with
relation $x=(n\bmod 2=1)$.  Reveal $n$ first and $x$ second.  The same rule
publishes both ordinary values when parity matches and publishes a failed $x$
when it does not.  No homogeneous universe or fabricated default member of
$\mathsf{Bool}$ or $\mathsf{Fin}(3)$ is required.

\subsection{Public chance before a completing reveal}

Suppose $y$ is public while its relation with the unrevealed $x$ is still
undecided, and a later public chance kernel reads $y$.  The draw is legitimately conditioned
on the already observed ordinary $y$.  If $x$ later closes the relation falsely,
$x$ fails; neither $y$ nor the chance result is rewritten.  Thus the semantics
is trace-sensitive.  A final-store interpretation which pretends the chance
had instead seen a null input is a different and generally false semantics.

\subsection{Private initial input}

An owner may know an initial private $x$ and choose $y$ satisfying a relation
before $x$ is publicly disclosed.  The initial secret and the public cell its
reveal creates must remain distinct.  If a dishonest earlier $y$ causes the later disclosure
to close the relation falsely, the publication of $x$ can fail while the ideal
private input remains $x$.  Honest automatic disclosure is an availability
assumption; null-vacuous guards do not establish it.

\section{Optional ordinary feasibility}

The all-failed run is consistent. A separate ordinary-feasibility witness
consists of an ordinary assignment $a_i:A_i$ satisfying every $R_g$, together
with bindings and openings that publish those same values.  Revealing the cells
in declared order then never fails: each guard is checked once, at which point
its support is entirely ordinary and its relation is true by the witness.

For a source compiler, the intended witness is the honest source execution.
The proof must show that every deferred relation is the reification of the
source guard in its exact choice-time view, that public inputs retain their
published values, and that private inputs later open to the same values the
owner used. This establishes an ordinary path, but is not required for
structural well-formedness or execution. In particular, a unary constant-false
guard is admitted: every ordinary publication is rejected and the resulting
failed publication discharges the lifted guard vacuously.

\section{Failure and strategic scope}

The revised source strategy space includes binding, reveal, and failure choices.
The intended compiler objective is therefore an exact causal mixture of source
deviations, not a generic premise that informed failure is utility-dominated.
Opponents remain unchanged and one admissible adaptive environment is fixed
across the comparison. Utility inequalities are needed only when a later target
adds actions, observations, timing power, or costs absent from the source, and
their strength must be justified separately.

A finite incentive example separating these points is not part of the checked
development. Whether informed failure is dominated is a property of a concrete
payoff structure, not of this semantics: the source language admits both, which
is why the compiler objective is an exact mixture law rather than a dominance
premise.

\section{Public implementation boundary}

The resolution rule has a direct smart-contract implementation when:
\begin{enumerate}
\item every guard has finite declared support and executable, deterministic
      typed code for $R_g$;
\item every support cell eventually receives an ordinary or failed public
      result;
\item a value opening is authenticated against the cell's immutable binding;
      and
\item the contract finalizes invalid and absent openings at the declared
      source reveal point, atomically checking the guards that reveal completes.
\end{enumerate}
The guards each reveal completes are known from the program text, so the
contract needs no per-cell status word and no unresolved-support counters: each
reveal carries its own fixed list of checks, evaluates their relations on the
already public values and the current proposal, and records failure for the
current cell if any is false.  Failed support makes the affected guards
vacuously satisfied, so no rollback or retraction is needed.

The semantic \lean{GuardCode.accepts} function is not itself a target-code
certificate. A compiler must provide publicly available expression code and
prove that its value inputs are exactly the declared support; a closure which
silently captures a secret is not an executable public validator. Ordinary
value encodings and guard evaluation must fit the target's resource limits.
For a finite program this requires finite storage and bounded local guard
evaluation, without changing the logical treatment of failure. Execution costs
and compensation for them remain part of the eventual strategic target model.

If a guard has an unresolved secret support member and no failed member, the
public evaluator cannot certify its nonvacuous truth.  A semantics demanding a
fully resolved terminal guard then needs disclosure, an explicit proof system,
an oracle, or another mechanism.  This is not a blanket runtime impossibility:
the guard needs no ordinary evaluation if some support member has already
failed and discharged it, and an application may explicitly permit an undecided
guard at its terminal boundary.  Likewise, an unreified semantic function or
unbounded evaluation cost requires an additional mechanism.  The finite model
should represent these cases honestly rather than calling private evaluation
public validation.

Target mempool contents, failed calls, retries, transaction ordering, gas, and
front-running are absent from this source model. The separate compiler proof must
couple outcome laws and causally translate deviations while retaining the
opponents and fixing the environment across the comparison. It must account
for the extra observations rather than assume they are functions of the source
view. The present consistency invariant supplies a useful
settlement safety lemma, but by itself proves neither that invalid and absent
traffic are strategically indistinguishable nor that native deviations admit
source-policy backtranslation.

\section{Publication invariant and compiler boundary}

The publication rule establishes:
\begin{quote}
Starting from a consistent public state, revealing one private cell by ordinary
publication when every guard that reveal completes accepts, and by failed
publication otherwise, preserves consistency, writes exactly one public cell,
and never retracts prior publication.  Under owned unrevealed support, every
rejecting guard which causes the fallback has a subject with the same owner as
the failed current publisher.
\end{quote}

Together with an optional ordinary witness, this shows that deferred relational
guards are internally consistent and can have failure-free executions. Separate
compiler obligations derive owned unrevealed support from source visibility,
prove exact honest source/publication correspondence, attribute each arbitrary
failure to a legal source alternative at the retained information prefix, and
transport the resulting outcome law. These obligations are discharged in the
source, graph, and pending-message proofs; none is assumed as an axiom of the
publication model.

\section{Source execution}
\label{sec:source}

A source state contains ordinary public inputs, the immutable private bindings
and published results of section \ref{sec:publications}, and the values of
chance draws. Both a binding and a publication of type $A$ are a
$\mathsf{Result}(A)$, so $\mathsf{ok}(\mathsf{none})$ for
$A=\mathsf{Option}(B)$ is a successful publication. Expressions eliminate
results explicitly; no ordinary default is supplied on failure.

Execution carries a registry of typed guards and, for each player, the ordered
history of its own binding and disclosure decisions. A player's policy is a
finite-distribution kernel on its public observations, its own private cells,
and this own-action history. The four source operations are:
\begin{itemize}
\item \emph{Commit:} draw a binding from the owner's policy and register the
written guard. No validity restriction is placed on the binding choice. The
guard is not evaluated here: its subject is not yet public.
\item \emph{Reveal:} draw a disclosure Boolean from the owner's policy. A
disclosure of an ordinary binding proposes its fixed value; withholding or an
unopenable binding proposes failure. Evaluate exactly the registered guards
this reveal completes -- those whose subject and code-read cells are all public
after it and were not all public before it -- against the proposal and the
earlier publications. Publish the proposal if every such guard accepts, and
failure otherwise, as a fresh public result alias. Remember the Boolean even
when both choices produce failure.
\item \emph{Sample:} draw an ordinary public value from the source's declared
kernel evaluated on the public environment. This law holds conditional on the
full execution prehistory, not only marginally. There is no strategic
disclosure choice for chance.
\item \emph{Return:} produce the detailed terminal state. Programmer-written
payoff expressions are a separate projection; a game analysis supplies utility.
\end{itemize}

The syntax tracks a finite set of unrevealed private resources. Every initial
private resource is unrevealed and accounted for; commitment adds one fresh
resource, reveal removes exactly one, and return requires the empty set.
Alongside it, the syntax carries a static map from each private cell to the
public cell its reveal creates, if any. Guards may read public inputs and their
author's own private cells, but not another player's bindings. All guard
support is static.

\paragraph{Static revelation.}
For any accounted program, the terminal map reveals every private cell. This is
a property of the syntax: no execution, policy, or chance outcome is involved,
and it is what makes every guard's completing reveal well defined.

\paragraph{Terminal safety.}
For any accounted initial state and any source policy profile, every outcome in
the support of complete execution decides each registered guard by its code:
either the publication of its subject or of a cell its code reads failed, or all
of them succeeded and the code holds on the published values. This includes
invalid bindings, withholding, and unsatisfiable ordinary guards.

The proof is an induction on source syntax carrying the invariant that every
guard whose components are all public already accepts. A reveal that completes
a guard publishes its proposal only when that guard accepts, and otherwise
publishes failure, which discharges the guard. A reveal that completes no guard
leaves every verdict unchanged, since accepted verdicts read only cells that
were already public.

\section{Checked interface}

\lean{Vegas.Foundation.Guard} contains the guard code and its verdict.
\lean{GuardCode.accepts} takes the subject's published result and results for
exactly the inputs its code reads; \lean{eq_accepts} characterizes it as the
unique such verdict which discharges on a failed input and otherwise evaluates
the code. \lean{Vegas.Source.Basic} adds the static revelation map
\lean{Revelations}, with \lean{Revelation.result} reading a published cell.

\lean{Vegas.Source} defines the typed source syntax, observation-local policies,
and exact finite-distribution execution above. \lean{SourceProgram.Initial}
packages initial resource accounting. Its principal safety theorems are:
\begin{itemize}
\item \lean{Initial.revealed}: every commitment is resolved by the end of the
syntax, with no execution quantifier;
\item \lean{Initial.terminal_guards_hold}: every retained guard is discharged by a
failed publication or holds on its published inputs.
\end{itemize}
\lean{Paper.lean} delegates directly to them. They do not assume successful
play or guard feasibility. \lean{VegasTests.SourceSemantics} exercises every
source constructor in one program with reverse disclosure, checking which
reveal completes the guard, its verdict there, and the full terminal-state law
under successful, rejecting, unopenable, and withholding policies.

\end{document}
